% Figura 2 de la postulación: índices de severidad (SV) y de peligro (H), con sus notas y niveles.
\documentclass[border=4pt]{standalone}
% Preámbulo común de las figuras 1 a 4 de la postulación (se carga con \input en cada figura).
% Los cortes y umbrales vienen de umbrales_calibrados.tex (lo escribe figuras_postulacion.py).
\usepackage{fontspec}
\setmainfont{Cambria}
\usepackage[table]{xcolor}
\usepackage{array,booktabs}
\usepackage{tikz}
\usetikzlibrary{arrows.meta,positioning,calc,fit,backgrounds,decorations.pathreplacing}
\definecolor{fondo}{HTML}{F2F1EC}
\definecolor{txt2}{HTML}{52514E}
% estados, ordenados por magnitud: E3 despeje < E4 tormenta < E5 peligro
\definecolor{cE0}{HTML}{D9D8D2}\definecolor{cE1}{HTML}{1BAF7A}\definecolor{cE2}{HTML}{EDA100}
\definecolor{cE3}{HTML}{2A78D6}\definecolor{cE4}{HTML}{E34948}\definecolor{cE5}{HTML}{4A3AA7}
\definecolor{cSV}{HTML}{2A78D6}\definecolor{cH}{HTML}{EB6834}\definecolor{cN}{HTML}{52514E}
% Generado por figuras_postulacion.py desde salidas/sistema/parametros.json. No editar a mano.
\newcommand{\cortesS}{\textbf{0{,}5} & 5 & \textbf{12{,}5} & 20 & \textbf{25}}
\newcommand{\cortesF}{0{,}05 & 0{,}25 & 0{,}50 & 0{,}75 & 0{,}95}
\newcommand{\cortesr}{0{,}5 & 1 & 2 & 3 & 4}
\newcommand{\cortesdur}{1 & 6 & 12 & 18 & 23}
\newcommand{\cortesG}{10 & 15 & 20 & 25 & 30}
\newcommand{\cortesR}{\textbf{0{,}5} & 2{,}5 & 5 & \textbf{10} & 15}
\newcommand{\cortesHNS}{10 & 20 & 30 & 50 & 80}
\newcommand{\cortesRS}{10 & 20 & 40 & 60 & 100}
\newcommand{\pesoS}{0{,}45}
\newcommand{\pesoF}{0{,}15}
\newcommand{\pesor}{0{,}10}
\newcommand{\pesodur}{0{,}10}
\newcommand{\pesoG}{0{,}10}
\newcommand{\pesoR}{0{,}10}
\newcommand{\SVrest}{1}
\newcommand{\SVparc}{2{,}5}
\newcommand{\SVcierre}{3{,}5}
\newcommand{\Hrest}{3}
\newcommand{\Hparc}{4}
\newcommand{\Hcierre}{5}
\newcommand{\Nrest}{5}
\newcommand{\Nparc}{30}
\newcommand{\Ncierre}{50}
\newcommand{\cdesp}{6}
\newcommand{\expAlto}{1}
\newcommand{\expMedio}{0{,}5}
\newcommand{\hnSobrecarga}{30}

\newcolumntype{L}[1]{>{\raggedright\arraybackslash}p{#1}}
\newcolumntype{C}{>{\centering\arraybackslash}p{0.62cm}}
% encabezado de notas con degradado del color del índice
\newcommand{\notas}[1]{\cellcolor{#1!12}1 & \cellcolor{#1!28}2 & \cellcolor{#1!45}3 & \cellcolor{#1!65}\textcolor{white}{4} & \cellcolor{#1!85}\textcolor{white}{5}}
\tikzset{
  tit/.style={anchor=north west, font=\small\bfseries, inner sep=0pt},
  caja/.style={draw=#1, line width=0.9pt, rounded corners=3pt, align=center, minimum width=4.2cm, minimum height=1.05cm, fill=white},
  dato/.style={draw=txt2, dashed, rounded corners=3pt, align=center, minimum width=4.5cm, minimum height=0.85cm},
  est/.style={rounded corners=2pt, align=center, minimum width=1.98cm, minimum height=1.05cm, font=\scriptsize, text=white, inner sep=1pt},
  niv/.style={rounded corners=2pt, align=center, minimum width=2.45cm, minimum height=0.62cm, font=\scriptsize, inner sep=1pt},
  llave/.style={decorate, decoration={brace, amplitude=3pt}, draw=txt2},
  flecha/.style={-{Stealth[length=2mm]}, thick, draw=txt2},
  eti/.style={font=\scriptsize, color=txt2, align=center}}

\newcolumntype{D}{>{\centering\arraybackslash}p{0.5cm}}
% fila de texto que ocupa todo el ancho de la tabla b) y corta en dos líneas si hace falta
\newcommand{\fila}[1]{\multicolumn{6}{>{\raggedright\arraybackslash}p{\dimexpr 6.8cm+25pt\relax}}{#1}}
% celda que ocupa las cinco columnas de notas: ajustes que se aplican a la nota de avalancha
\newcommand{\ajuste}[1]{\multicolumn{5}{>{\centering\arraybackslash}p{\dimexpr 2.5cm+20pt\relax}}{#1}}
\begin{document}
\begin{tikzpicture}[font=\footnotesize]
% a) Severidad
\node[tit] at (0,0) {a) Severidad de la tormenta, $SV$ (0 a 5)};
\node[anchor=north west, inner sep=0pt] (tb) at (0,-0.45) {\scriptsize\setlength{\tabcolsep}{2.5pt}\renewcommand{\arraystretch}{1.18}%
\begin{tabular}{L{3.55cm}DDDDD>{\centering\arraybackslash}p{0.62cm}}
\toprule
Variable medida en cada bloque de 6 horas & \notas{cSV} & peso \\
\midrule
Nieve acumulada sobre el camino (cm) & \cortesS & \pesoS \\
Fracción del camino cubierta de nieve & \cortesF & \pesoF \\
Tasa de nevada (cm/h) & \cortesr & \pesor \\
Horas que nevó durante las últimas 24 horas & \cortesdur & \pesodur \\
Ráfaga máxima de viento (m/s) & \cortesG & \pesoG \\
Lluvia sobre el manto de nieve (mm) & \cortesR & \pesoR \\
\midrule
\multicolumn{7}{l}{$SV$ = máx\{piso; $\sum$ peso $\times$ nota\}} \\
\bottomrule
\end{tabular}};
\path (tb.south west) -- (tb.south east)
  node[pos=0.17, below=0.55cm, niv, fill=cE1, text=white] {restricción\\$SV \ge \SVrest$}
  node[pos=0.5, below=0.55cm, niv, fill=cE2] {cierre parcial\\$SV \ge \SVparc$}
  node[pos=0.83, below=0.55cm, niv, fill=cE4, text=white] {cerrado por tormenta\\$SV \ge \SVcierre$};

% b) Peligro
\node[tit] at ($(tb.north east)+(0.6,0.45)$) (tbt) {b) Peligro de avalancha y remoción, $H$ (0 a 5)};
\node[anchor=north west, inner sep=0pt] (tc) at ($(tb.north east)+(0.6,0)$) {\scriptsize\setlength{\tabcolsep}{2.5pt}\renewcommand{\arraystretch}{1.18}%
\begin{tabular}{L{4.3cm}DDDDD}
\toprule
Variable & \notas{cH} \\
\midrule
\fila{\cellcolor{fondo}\textit{Avalancha}} \\
Nieve nueva acumulada durante las últimas 72 horas (cm) & \cortesHNS \\[2pt]
Viento que transporta la nieve nueva & \ajuste{+1 a la nota, con \hnSobrecarga\ cm o más} \\[2pt]
Isoterma que sube 500 m en las últimas 24 horas, o 5 mm de lluvia sobre el manto & \ajuste{+1 a la nota, con \hnSobrecarga\ cm o más} \\[2pt]
Exposición del tramo a avalanchas & \ajuste{$\times$ \expAlto\ en el tramo alto\newline$\times$ \expMedio\ en el tramo medio} \\[2pt]
\fila{\cellcolor{fondo}\textit{Remoción en masa}, solo con la isoterma sobre la cota media del tramo} \\
Lluvia acumulada durante las últimas 72 horas (mm) & \cortesRS \\
\midrule
\fila{$H$ = máx\{avalancha; remoción en masa\}} \\
\bottomrule
\end{tabular}};
\path (tc.south west |- tb.south west) -- (tc.south east |- tb.south east)
  node[pos=0.17, below=0.55cm, niv, fill=cE1, text=white] {restricción\\$H \ge \Hrest$}
  node[pos=0.5, below=0.55cm, niv, fill=cE2] {cierre parcial\\$H \ge \Hparc$}
  node[pos=0.83, below=0.55cm, niv, fill=cE5, text=white] {cerrado por peligro\\$H = \Hcierre$};
\end{tikzpicture}
\end{document}
